%%%%%%%%%%%%%%%%%%%%%%%%%%%%%%%%%%%%%%%%%
% Professional Formal Letter
% LaTeX Template
% Version 2.0 (12/2/17)
%
% This template originates from:
% http://www.LaTeXTemplates.com
%
% Authors:
% Brian Moses
% Vel (vel@LaTeXTemplates.com)
%
% License:
% CC BY-NC-SA 3.0 (http://creativecommons.org/licenses/by-nc-sa/3.0/)
%
%%%%%%%%%%%%%%%%%%%%%%%%%%%%%%%%%%%%%%%%%

%----------------------------------------------------------------------------------------
%	PACKAGES AND OTHER DOCUMENT CONFIGURATIONS
%----------------------------------------------------------------------------------------

\documentclass[11pt, a4paper]{letter} % Set the font size (10pt, 11pt and 12pt) and paper size (letterpaper, a4paper, etc)

%%%%%%%%%%%%%%%%%%%%%%%%%%%%%%%%%%%%%%%%%
% Professional Formal Letter
% Structure Specification File
% Version 1.0 (12/2/17)
%
% This file originates from:
% http://www.LaTeXTemplates.com
%
% Authors:
% Brian Moses
% Vel (vel@LaTeXTemplates.com)
%
% License:
% CC BY-NC-SA 3.0 (http://creativecommons.org/licenses/by-nc-sa/3.0/)
%
%%%%%%%%%%%%%%%%%%%%%%%%%%%%%%%%%%%%%%%%%

%----------------------------------------------------------------------------------------
%	PACKAGES AND OTHER DOCUMENT CONFIGURATIONS
%----------------------------------------------------------------------------------------

\usepackage{graphicx} % Required for including pictures

\def\magyarOptions{hyphenation=huhyphn}
\usepackage{ae,aecompl}
\usepackage[T1]{fontenc} % Output font encoding for international characters
\usepackage[utf8]{inputenc} % Required for inputting international characters
\usepackage[hungarian]{babel}
\usepackage{indentfirst}
\usepackage{enumerate}
\usepackage{setspace}
\setstretch{1.2}

%\usepackage{gfsdidot} % Use the GFS Didot font: http://www.tug.dk/FontCatalogue/gfsdidot/
%\usepackage{microtype} % Improves typography

%\pagestyle{empty} % Suppress headers and footers

\setlength\parindent{1cm} % Paragraph indentation

% Create a new command for the horizontal rule in the document which allows thickness specification
\makeatletter
\newcommand{\vhrulefill}[1]{\leavevmode\leaders\hrule\@height#1\hfill \kern\z@}
\makeatother

%----------------------------------------------------------------------------------------
%	DOCUMENT MARGINS
%----------------------------------------------------------------------------------------

\usepackage{geometry} % Required for adjusting page dimensions

\geometry{
	top=2cm, % Top margin
	bottom=2.5cm, % Bottom margin
	left=3cm, % Left margin
	right=3cm, % Right margin
	%showframe, % Uncomment to show how the type block is set on the page
}

%----------------------------------------------------------------------------------------
%	DEFINE CUSTOM COMMANDS
%----------------------------------------------------------------------------------------

\newcommand{\logo}[1]{\renewcommand{\logo}{#1}}

\newcommand{\Who}[1]{\renewcommand{\Who}{#1}}
\newcommand{\Title}[1]{\renewcommand{\Title}{#1}}

\newcommand{\headerlineone}[1]{\renewcommand{\headerlineone}{#1}}
\newcommand{\headerlinetwo}[1]{\renewcommand{\headerlinetwo}{#1}}

\newcommand{\authordetails}[1]{\renewcommand{\authordetails}{#1}}

%----------------------------------------------------------------------------------------
%	AUTHOR DETAILS STRUCTURE
%----------------------------------------------------------------------------------------

\newcommand{\authordetailsblock}{
	\hspace{\fill} % Move the author details to the far right
	\parbox[t]{0.48\textwidth}{ % Box holding the author details; width value specifies where it starts and ends, increase to move details left
		\footnotesize % Use a smaller font size for the details
		\Who\\ % Author name
		\textit{\authordetails} % The author details text, all italicised
	}
}

%----------------------------------------------------------------------------------------
%	HEADER STRUCTURE
%----------------------------------------------------------------------------------------

\address{
	\includegraphics[width=1in]{\logo} % Include the logo of author institution
	\hspace{0.82\textwidth} % Position of the institution logo, increase to move left, decrease to move right
	\vskip -0.1\textheight~\\ % Position of the large header text in relation to the institution logo, increase to move down, decrease to move up
	\Large\hspace{0.2\textwidth}\headerlineone\hfill ~\\[0.006\textheight] % First line of institution name, adjust hspace if your logo is wide
	\hspace{0.2\textwidth}\headerlinetwo\hfill \normalsize % Second line of institution name, adjust hspace if your logo is wide
	\makebox[0ex][r]{\textbf{\Who\Title}}\hspace{0.01\textwidth} % Print author name and title with a little whitespace to the right
	~\\[-0.01\textheight] % Reduce the whitespace above the horizontal rule
	\hspace{0.2\textwidth}\vhrulefill{1pt} \\ % Horizontal rule, adjust hspace if your logo is wide and \vhrulefill for the thickness of the rule
	\authordetailsblock % Include the letter author's details on the right side of the page under the horizontal rule
	\hspace{-0.25\textwidth} % Horizontal position of the author details block, increase to move left, decrease to move right
	\vspace{-0.1\textheight} % Move the date and letter content up for a more compact look
}

%----------------------------------------------------------------------------------------
%	COMPOSE THE ENTIRE HEADER
%----------------------------------------------------------------------------------------

\renewcommand{\opening}[1]{
	{\centering\fromaddress\vspace{0.05\textheight} \\ % Print the header and from address here, add whitespace to move date down
	\hspace*{\longindentation}\today\hspace*{\fill}\par} % Print today's date, remove \today to not display it
	{\raggedright \toname \\ \toaddress \par} % Print the to name and address
	\vspace{1cm} % White space after the to address
	\noindent #1 % Print the opening line
}

%----------------------------------------------------------------------------------------
%	SIGNATURE STRUCTURE
%----------------------------------------------------------------------------------------

\signature{\Who\Title} % The signature is a combination of the author's name and title

\renewcommand{\closing}[1]{
	\vspace{2.5mm} % Some whitespace after the letter content and before the signature
	\noindent % Stop paragraph indentation
	\hspace*{\longindentation} % Move the signature right to the value of \longindentation
	\parbox{\indentedwidth}{
		\raggedright
		#1 % Print the signature text
		\vskip 1.65cm % Whitespace between the closing text and author's name for a physical signature
		\fromsig % Prints the value of \signature{}, i.e. author name and title
	}
}
 % Include the file that specifies the document structure

%\longindentation=0pt % Un-commenting this line will push the closing "Sincerely," and date to the left of the page

%----------------------------------------------------------------------------------------
%	YOUR INFORMATION
%----------------------------------------------------------------------------------------

\Who{Kiss András} % Your name

\Title{, PhD} % Your title, leave blank for no title

\authordetails{
	Ált. és Fiz. Kém. Tsz.\\ % Your department/institution
	Ifjúság útja 6.\\ % Your address
	Pécs, 7624\\ % Your city, zip code, country, etc
	akiss@gamma.ttk.pte.hu\\ % Your email address
	Belső mellék: 61021\\ % Your phone number
}

%----------------------------------------------------------------------------------------
%	HEADER CONTENTS
%----------------------------------------------------------------------------------------

\logo{pte_logo.eps} % Logo filename, your logo should have square dimensions (i.e. roughly the same width and height), if it does not, you will need to adjust spacing within the HEADER STRUCTURE block in structure.tex (read the comments carefully!)

\headerlineone{PÉCSI} % Top header line, leave blank if you only want the bottom line

\headerlinetwo{TUDOMÁNYEGYETEM} % Bottom header line

%----------------------------------------------------------------------------------------

\begin{document}

%----------------------------------------------------------------------------------------
%	TO ADDRESS
%----------------------------------------------------------------------------------------

\begin{letter}{
%	Prof. Jones\\
%	Mathematics Search Committee\\
%	Department of Mathematics\\
%	University of California\\
%	Berkeley, California 12345
}

%----------------------------------------------------------------------------------------
%	LETTER CONTENT
%----------------------------------------------------------------------------------------

\opening{
\begin{center}Témavezetői vélemény Szili Szilárd \emph{,,A Belouszov--Zsabotyinszkij reakció pásztázó elektrokémiai mikroszkópos térképezése''} című szakdolgozatáról
\end{center}
}

%\flushleft

Szilárd 2017-ben kezdett témavezetésem mellett foglalkozni az elektrokémiai mikroszkóppal. Tanszékünk elektrokémiai kutatócsoportja főleg mikroskálán megfigyelhető jelenségek elektroanalitikai vizsgálatával foglalkozik, így Szilárd az ehhez szükséges mikroelektródok készítésével és karakterizálásával kezdte a munkát. A 2017 június elejétől augusztus végéig terjedő nyári gyakorlata során ezen kívül megismerkedett az alapvető elektroanalitikai módszerekkel is. Szilárdnak eredetileg egy egyszerű, már viszonylag jól vizsgált témát jelöltem ki. Az antimon mikroelektróddal történő pH térképezés további lehetőségeit és akadályait kezdte vizsgálni. A munkát rendkívül nagy fegyelemmel és szorgalommal végezte. Nyári gyakorlata során nyílvánvalóvá vált számomra, hogy Szilárd ennél jóval többre lesz képes, és a választottnál érdekesebb, egyben nagyobb kihívást nyújtó témát érdemel. Ezért megkérdeztem Szilárdtól, hogy szeretne-e egy olyan kutatáshoz csatlakozni, mely a pásztázó elektrokémiai mikroszkópot egy teljesen új módon használja oldatfázis felületi pásztázására. Szilárd nagy lelkesedéssel kezdett dolgozni az új problémán. Az irodalomban erre még nem volt példa, a téma nekem is új volt. Mivel ritka jelenség, hogy homogén oldatfázisban mintázat alakul ki, a például szolgáló rendszer kiválasztása könnyű volt. Az kémiai oszcilláló jelenségek közül a már nagyon jól jellemzett Belouszov-Zsabotyinszkij oszcilláló reakciót (BZ-reakció) választottam. Szilárd a BZ-reakcióban kialakuló redox-potenciál mintázatot térképezte pásztázó elektrokémiai mikroszkóppal. 2017 szeptembere és 2018 áprilisa között végezte a kísérletes munkát. Az erről, körülbelül három hónap leforgása alatt írt szakdolgozata a hivatkozásjegyzékkel együtt 32 számozott oldalból áll, így megfelel a terjedelemmel szemben támasztott követelményeknek. A dolgozat arányos felépítésű, nagy része körülbelül egyenlően van felosztva a három nagy fejezetre (Irodalom, Módszerek, Eredmények). A Bevezetés kellő távolságból közelít rá a választott témára. A Célkitűzések fejezet világosan, pontokba szedve írja le a munka célját. A dolgozat irodalomjegyzéke igényes, egységes formátumú, hibát csak elvétve találni (, pl. [27]). Dícséretes, hogy a dolgozat tizenöt ábrája közül tizennégy Szilárd sajátkészítésű ábrája. A dolgozat LaTeX szövegformázó rendszerrel készült, éppen ezért tipográfiailag helyes, és esztétikus. Néhány ábra azonban rossz betűméretet használ, például a 4.6 ábra betűi túl nagyok, az 5.2 és 4.5 betűi túl kicsik. A LaTeX előnyei mellett egyik nagy hátránya az automatikus helyesírás ellenőrzés hiánya, és ez sajnos meg is látszik a dolgozaton. Viszonylag sok az elgépelés. A dolgozat végső verziójának produkálása és nyomtatása előtt egy alaposabb ellenőrzésre lett volna szükség. Így valószínűleg elkerülhetőek lettek volna a 3.3 és 3.4 egyenletben lévő sztöchiometriai hibák is. A felsorolt hibák mind formai jellegűek, az egyenletekben vétett hibák pedig természetesen csak pillanatnyi figyelmetlenség eredményei. A dolgozat nyelvezete jó, de helyenként nehézkes a fogalmazás. Szilárdnak mindenképpen javítania kell az írásbeli kifejezőképességén.

A kitűzött célokat végül sikerült elérni. Ezt nagy teljesítménynek tartom Szilárd részéről, hiszen nem csupán egy már jól ismert témának egy kis részletén dolgozott, hanem a kutatócsoportunkban használt technika egy teljesen új alkalmazását dolgozta ki témavezetésem mellett. Meggyőződtem róla, hogy Szilárd a kísérletes munka és a szakdolgozat írása során alkalmazott módszerek minden részletét érti, azokat magas szinten elsajátította. A dolgozat messzemenően kielégíti a BSc. szakdolgozattal szemben támasztott követelményeket. A védésre a következő kérdéseket teszem fel:

\begin{enumerate}
\item A kísérleti paraméterek milyen változtatása szükséges élő rendszerekben megjelenő kémiai hullámok, kémiai mintázatok térképezéséhez?
\item Vizsgálta-e a BZ-reakciót kisebb skálán, esetleg kevesebb, vagy több dimenzióban? Ha igen, milyen tapasztalatai vannak?
\end{enumerate}

\begin{center}
\textbf{Sikeres védés esetén javaslom a jeles érdemjegy odaítélését.}
\end{center}

Kiss András
Egyetemi tanársegéd
Általános és Fizikai Kémia Tanszék
Kémiai Intézet
Pécsi Tudományegyetem


2018, június 16
Homburg (Saar)
Németország

%\closing{Sincerely,}

%----------------------------------------------------------------------------------------
%	OPTIONAL FOOTNOTE
%----------------------------------------------------------------------------------------

% Uncomment the 4 lines below to print a footnote with custom text
%\def\thefootnote{Témavezetői vélemény}
%\def\footnoterule{\hrule}
%\footnotetext{\hspace*{\fill}{\footnotesize\em Footnote text}}
%\def\thefootnote{\arabic{footnote}}

%----------------------------------------------------------------------------------------

\end{letter}

\end{document}
